\section{Zusammenfassung der Ergebnisse}

Die vorliegende Arbeit hatte zum Ziel, eine strukturierte Risikomanagement-Strategie 
für Data-Science-Projekte zu entwickeln, die als Erweiterung des DASC-PM konzipiert 
ist. Ausgehend von einer Analyse zentraler Risikotypen und etablierter Methoden der 
Risikobeurteilung, -minimierung und des -monitorings wurde das Konzept der 
\textit{Risk Gates} als phasenübergreifendes Steuerungsinstrument hergeleitet und 
definiert. Die entwickelte \textit{RM-Extension} integriert dieses Konzept in die 
Hauptphasen des DASC-PM und operationalisiert es über ein dreistufiges Ampelsystem, 
das auf messbaren Key Risk Indicators (KRI) basiert.

Die Anwendung auf zwei Fallstudien -- das fiktive Kreditprüfungsprojekt \hyperref[sec:case_study_1]{SafeLend}
sowie die reale Personalisierungsinitiative \hyperref[sec:case_study_2]{Jumbo} -- hat gezeigt, dass die RM-Extension auf unterschiedliche Projektkontexte übertragbar ist.
In beiden Fällen ermöglichte das strukturierte Gate-Verfahren eine systematische Priorisierung von Risiken, eine klare Entscheidungsgrundlage für Phasenübergänge sowie eine nachvollziehbare Dokumentation des Risikomanagementprozesses.
Die Carry-Forward-Logik des Gelb-Status hat sich dabei als zentrales Instrument erwiesen, um erhöhte, jedoch noch nicht kritische Risiken über mehrere Phasen hinweg kontinuierlich im Blick zu behalten, ohne den Projektfortschritt unnötig zu blockieren.

